\scanpage{52}
\mspage{III}

\uline{Lemme 2.} Soit $X\longrightarrow Y$ un morphisme de type fini, avec $Y$ loc. noeth. et univ. japonais. Alors pour tt faisceau d'ensembles constructible $F$ sur $X$, $f_*(F)$ est constructible.

Standart

\uline{Lemme 3} $X$ schéma connexe \deleted{\unsure{loc.}} noeth, \inserted{univ. japonais}, $(X_i)_{i\in I}$ une partition finie en parties loc fermées, $\unsure{G|C}$ un gpe fini. \deleted{Supposons que pour tout $i$, $H^1(X_i,G)$ soit fini.} \inserted{Considérons l'\deleted{application}}
\deleted{Alors \unsure{pour tt}}
\[
H^1(X,G)\longrightarrow\prod_i H^1(X_i,G)
\]

Cette application est à \uline{fibres finies}. En particulier, si pour tt $i$, $H^1(X_i,G)$ est fini, Kif Kif pour $H^1(X,G)$.

Théorie de recollement des faisceaux étales, en utilisant lemme 2. et le fait suivant :
\inserted{\uline{Lemme 4}} Si $F,G$ sont des faisceaux \inserted{d'ens} étales \unsure{constr.} sur $X$ noethérien, alors $\operatorname{Hom}(F,G)$ est un ens fini.
\margin{$\big|$}

\deleted{\uline{Proposition}} \deleted{\uline{\unsure{Corollaire|V$^{\mathrm{o}}$ Main Proposition}}} \inserted{\uline{Proposition}} $S$ \deleted{\unsure{$X$}} schéma noeth. univ. japonais, $G$ un gpe fini. Conditions équivalentes :
% archive sheet number: 51
